\documentclass[11pt]{article}
\input{../../shared/project_style.tex}
\rhead{Basics 12 Textbook Draft}

\begin{document}
\DocTitle{Basics 12: Maxwell Equations}{II - Electromagnetism foundations}{Local and integral laws, conservation of charge, electromagnetic waves, and plasma limits}

\begin{overviewbox}
\textbf{Textbook draft, version 1.0.} Maxwell's equations organize classical electromagnetism into four coupled field laws. This chapter derives their integral forms, obtains charge conservation and vacuum electromagnetic waves, and then identifies the low-frequency approximations used in magnetohydrodynamics. This chapter is designed for a reader with no prior plasma-physics training. It distinguishes exact identities from physical approximations and connects the central equations to later simulation modules.
\end{overviewbox}

\ProjectTOC

\section{Learning objectives}
After completing this note, the reader should be able to:
\begin{itemize}[leftmargin=1.6em]
\item State Maxwell's equations in differential and integral form.
\item Derive the continuity equation for electric charge.
\item Derive the electromagnetic wave equation in vacuum.
\item Explain the displacement-current term.
\item Distinguish exact field constraints from model-dependent plasma approximations.
\item Identify numerical consequences of divergence constraints.
\end{itemize}

\section{Prerequisites and reading strategy}
\begin{itemize}[leftmargin=1.6em]
\item Basics 2, 3, and 11.
\end{itemize}
On a first reading, emphasize physical meaning and dimensional checks. On a second reading, reproduce the derivations. Attempt the computational laboratory only after the limiting cases are understood.

\section{The four differential equations}

In SI units,
\begin{align}
\nabla\cdot\mathbf E &= \frac{\rho_q}{\epsilon_0},\\
\nabla\cdot\mathbf B &= 0,\\
\nabla\times\mathbf E &= -\frac{\partial\mathbf B}{\partial t},\\
\nabla\times\mathbf B &= \mu_0\mathbf J+\mu_0\epsilon_0\frac{\partial\mathbf E}{\partial t}.
\end{align}
Gauss's electric law relates charge to electric flux. Gauss's magnetic law states that the net magnetic flux through a closed surface vanishes. Faraday's law describes induction. Ampere-Maxwell law connects magnetic circulation to conduction current and changing electric flux.


\section{Integral forms from the divergence and Stokes theorems}

Integrating the divergence equations over a volume $V$ gives
\begin{align}
\oint_{\partial V}\mathbf E\cdot d\mathbf A &= \frac{Q_{\rm enc}}{\epsilon_0},\\
\oint_{\partial V}\mathbf B\cdot d\mathbf A &=0.
\end{align}
Integrating the curl equations over an oriented surface $S$ gives
\begin{align}
\oint_{\partial S}\mathbf E\cdot d\boldsymbol\ell
&=-\frac{d}{dt}\int_S\mathbf B\cdot d\mathbf A,\\
\oint_{\partial S}\mathbf B\cdot d\boldsymbol\ell
&=\mu_0 I_{\rm enc}+\mu_0\epsilon_0\frac{d}{dt}\int_S\mathbf E\cdot d\mathbf A.
\end{align}
The integral forms expose global flux and circulation balances; the differential forms describe their local density.


\section{Why displacement current is required}

Taking the divergence of Ampere-Maxwell law and using $\nabla\cdot(\nabla\times\mathbf B)=0$ yields
\begin{equation}
0=\mu_0\nabla\cdot\mathbf J+\mu_0\epsilon_0
\frac{\partial}{\partial t}\left(\nabla\cdot\mathbf E\right).
\end{equation}
Gauss's law then gives
\begin{equation}
\frac{\partial\rho_q}{\partial t}+\nabla\cdot\mathbf J=0.
\end{equation}
Without the displacement-current term, a time-varying charge density would contradict local charge conservation. Maxwell's correction is therefore not an optional decoration.


\section{Electromagnetic waves in vacuum}

Set $\rho_q=0$ and $\mathbf J=0$. Taking the curl of Faraday's law gives
\begin{equation}
\nabla\times\nabla\times\mathbf E
=-\frac{\partial}{\partial t}\left(\nabla\times\mathbf B\right).
\end{equation}
Using $\nabla\times\nabla\times\mathbf E=\nabla(\nabla\cdot\mathbf E)-\nabla^2\mathbf E$ and $\nabla\cdot\mathbf E=0$,
\begin{equation}
\nabla^2\mathbf E-\mu_0\epsilon_0\frac{\partial^2\mathbf E}{\partial t^2}=0.
\end{equation}
The wave speed is
\begin{equation}
c=\frac{1}{\sqrt{\mu_0\epsilon_0}}.
\end{equation}
An identical equation follows for $\mathbf B$. Plane waves are transverse and satisfy $B=E/c$.


\section{Potentials and gauge freedom}

The identities
\begin{equation}
\mathbf B=\nabla\times\mathbf A,\qquad
\mathbf E=-\nabla\phi-\frac{\partial\mathbf A}{\partial t}
\end{equation}
automatically satisfy the two homogeneous Maxwell equations. The transformation
\begin{equation}
\mathbf A\rightarrow\mathbf A+\nabla\chi,\qquad
\phi\rightarrow\phi-\frac{\partial\chi}{\partial t}
\end{equation}
leaves $\mathbf E$ and $\mathbf B$ unchanged. Gauge freedom changes the mathematical representation, not measurable fields.


\section{Low-frequency plasma limits}

MHD typically assumes characteristic speeds much smaller than $c$ and frequencies low enough that
\begin{equation}
\epsilon_0\left|\frac{\partial\mathbf E}{\partial t}\right|\ll |\mathbf J|.
\end{equation}
Ampere-Maxwell law then becomes approximately
\begin{equation}
\nabla\times\mathbf B\simeq\mu_0\mathbf J.
\end{equation}
Quasineutrality further assumes $|\rho_q|\ll en$, but it does not imply $\mathbf E=0$. These approximations remove light waves and plasma-frequency charge oscillations from the resolved dynamics. Their domains of validity must be stated.


\section{Constraint preservation in numerical models}

The divergence equations act as constraints on admissible initial data. Faraday's law preserves $\nabla\cdot\mathbf B=0$ analytically because
\begin{equation}
\frac{\partial}{\partial t}(\nabla\cdot\mathbf B)
=-\nabla\cdot(\nabla\times\mathbf E)=0.
\end{equation}
A discrete scheme may violate the identity. Constrained transport, divergence cleaning, or compatible finite elements are therefore used to control numerical magnetic monopoles.


\section{Computational laboratory}
Discretize a one-dimensional vacuum electromagnetic wave and verify $c=1/\sqrt{\mu_0\epsilon_0}$. Independently evaluate differential and integral forms of Gauss's law for a smooth charge distribution. Measure discrete charge-conservation and divergence errors.

\section{Summary}
\begin{itemize}[leftmargin=1.6em]
\item Maxwell's equations couple charges, currents, and time-dependent fields.
\item The displacement current guarantees compatibility with charge conservation.
\item Vacuum electromagnetic waves follow directly from the curl equations.
\item MHD is a controlled low-frequency reduction, not a replacement for Maxwell theory.
\item Divergence constraints require deliberate preservation in numerical simulations.
\end{itemize}

\SymbolsAcronymsSection
\begin{symtable}
$\rho_q$ & charge density & Source in Gauss's electric law. \\
$\mathbf J$ & current density & Charge flow per area. \\
$c$ & speed of light & $1/\sqrt{\mu_0\epsilon_0}$. \\
$\mathbf A$ & vector potential & Satisfies $\mathbf B=\nabla\times\mathbf A$. \\
$\chi$ & gauge function & Generates equivalent potentials. \\
$Q_{\rm enc},I_{\rm enc}$ & enclosed source measures & Volume charge and surface-crossing current. \\
\end{symtable}

\section{Exercises}
\begin{enumerate}[leftmargin=1.8em]
\item Derive charge conservation from Ampere-Maxwell and Gauss laws.
\item Derive the vacuum magnetic-field wave equation.
\item Show that a plane electromagnetic wave is transverse.
\item Explain why quasineutrality does not eliminate electric fields.
\item Estimate when displacement current is negligible relative to conduction current.
\end{enumerate}

\section{References}
{\small
\begin{thebibliography}{99}
\bibitem{ref1} G. B. Arfken, H. J. Weber, and F. E. Harris, \emph{Mathematical Methods for Physicists}, 7th ed., Academic Press (2013).
\bibitem{ref2} G. Strang, \emph{Linear Algebra and Its Applications}, 4th ed., Brooks/Cole (2006).
\bibitem{ref3} W. A. Strauss, \emph{Partial Differential Equations: An Introduction}, 2nd ed., Wiley (2007).
\bibitem{ref4} D. J. Griffiths, \emph{Introduction to Electrodynamics}, 4th ed., Cambridge University Press (2017).
\bibitem{ref5} J. D. Jackson, \emph{Classical Electrodynamics}, 3rd ed., Wiley (1998).
\bibitem{ref6} F. F. Chen, \emph{Introduction to Plasma Physics and Controlled Fusion}, 3rd ed., Springer (2016).
\bibitem{ref7} R. J. Goldston and P. H. Rutherford, \emph{Introduction to Plasma Physics}, CRC Press (1995).
\bibitem{ref8} P. M. Bellan, \emph{Fundamentals of Plasma Physics}, Cambridge University Press (2006).
\bibitem{ref9} J. P. Freidberg, \emph{Ideal MHD}, Cambridge University Press (2014).
\bibitem{ref10} J. P. Goedbloed, R. Keppens, and S. Poedts, \emph{Advanced Magnetohydrodynamics}, Cambridge University Press (2010).
\bibitem{ref11} H. Goedbloed and S. Poedts, \emph{Principles of Magnetohydrodynamics}, Cambridge University Press (2004).
\end{thebibliography}
}

\end{document}
